\documentclass[10pt,a4paper]{amsart}
\usepackage[margin=19mm]{geometry}
\usepackage{amsmath,amssymb,mathtools}
\usepackage[scr=rsfs,cal=euler]{mathalfa}
\usepackage{xfrac}
\usepackage[unicode,hidelinks,bookmarks=false]{hyperref}
\newcommand{\ie}{i.e.\ }
\newcommand{\cf}{cf.\ }
\newcommand{\resp}{respectively\ }
\newcommand{\Subsection}{\subsection}
\providecommand{\nobreakdash}{\nobreak\hbox{-}}
\setlength{\emergencystretch}{2em}
\multlinegap0pt
\usepackage{bm}
\usepackage{bbm}
\usepackage{dsfont}
\usepackage{xspace}
\usepackage{cancel}
\usepackage{mathtools}
\usepackage{amsthm}
\makeatletter
\@ifundefined{theorem}{\newtheorem{theorem}{Theorem}}{}
\@ifundefined{lemma}{\newtheorem{lemma}[theorem]{Lemma}}{}
\@ifundefined{proposition}{\newtheorem{proposition}[theorem]{Proposition}}{}
\makeatother
\title[Harnack parts for 5-by-5 truncated shift with numerical radi]{Harnack parts for 5-by-5 truncated shift with numerical radius one}
\author{Mohammed Benharrat}
\date{Source: arXiv:2604.03670v1 (CC BY 4.0)}
\begin{document}
\maketitle

\noindent\emph{Source and licence.} Harnack parts for 5-by-5 truncated shift with numerical radius one. Version arXiv:2604.03670v1 (CC BY 4.0), distributed under the Creative Commons Attribution 4.0 International licence, as recorded in the arXiv licence metadata for this exact version and shown on the abstract page https://arxiv.org/abs/2604.03670v1. Full rights notice: licenses/arxiv-2604.03670-CC-BY-4.0.txt.\par\smallskip

\noindent\textit{Editorial scope.} Counted unit: the Proposition whose source label is \texttt{Prop:case2} in the article (source0.tex lines 624--626, source label \texttt{Prop:case2}), delivered together with the article's own complete original proof of it (source0.tex lines 675--828, one \texttt{proof} environment of 154 lines). No mathematics is written, repaired or re-derived: statement and proof are the article's own text line for line. The proof is detached from the statement in the source: the article states the result at lines 624--626 and prints its proof at lines 675--828, and the proof environment's own header names the statement, so the association is the article's own. Required context retained uncounted: form of T (lines 221--234); Harnack\_part\_S\_5 (lines 313--332); eqfour (lines 368--370); lemma:Rn\_unified (lines 630--652); c1 (lines 648--651); c2 (lines 648--651). Every definition, display and label of those lines is retained unchanged. Omitted from this package and not redistributed: 1--220, 235--312, 333--367, 371--623, 627--629, 652--674, 829--1009. Nothing inside a retained excerpt is omitted or altered: every retained body is delivered exactly as the article's lines are written, so no omission receipt of any kind accompanies this package. Citations: the retained excerpts carry no citation command at all, so no bibliography entry of the article is transcribed and no reference is redistributed. Reference adaptation: none was needed and none was made; every reference inside the delivered unit resolves inside it, so no unresolved reference or citation is produced. Printed slips kept literal: no printed word, symbol, hypothesis or number of the counted statement or of its proof was altered. Layout and class: the article's own class is \texttt{amsart} with packages including amsmath, amsthm, amscd, amsfonts, amssymb, graphicx, hyperref, color, inputenc, fontenc, enumitem, mathtools, empheq; the wrapper is the lane's pinned \texttt{amsart} editorial wrapper at 10pt on A4, which restates the theorem-like environments the retained text needs and adds the article's own macro definitions that the retained text needs (none), with extra wrapper packages (bm, bbm, dsfont, xspace, cancel, mathtools). The delivered numbering is the unit's own. Assets: no retained span contains a figure environment, a graphics-inclusion command, a PDF-page inclusion, a table or a TikZ picture; no image file is copied into this package, the package declares no asset dependency and no image file is redistributed with it. Licence: version arXiv:2604.03670v1 of this article is distributed under the Creative Commons Attribution 4.0 International licence, as recorded in the arXiv licence metadata for this exact version and shown on the abstract page https://arxiv.org/abs/2604.03670v1. A full rights notice is delivered at licenses/arxiv-2604.03670-CC-BY-4.0.txt. Characters: every retained excerpt is pure ASCII, so the delivered engine, XeLaTeX with its default Latin Modern text font, reports no missing character.
\begin{equation}\label{form of T}
	T = \begin{bmatrix}
		0 & a_1 & a_2 & a_3 & 0 \\
		0 & r_{11} & r_{12} & r_{13} & b_1 \\
		0 & r_{21} & r_{22} & r_{23} & b_2 \\
		0 & r_{31} & r_{32} & r_{33} & b_3 \\
		0 & 0 & 0 & 0 & 0 \\
	\end{bmatrix} 
	= \begin{bmatrix}
		0 & A & 0 \\
		0 & R & B \\
		0 & 0 & 0 \\
	\end{bmatrix},
\end{equation} 

\begin{proposition}\label{Harnack_part_S_5} 
	Assume that the spectrum of $R$ is $\sigma(R)=\{ 0, \lambda \}$, with $\lambda \in \mathbb{D}$, and that the eigenvalue $0$ has algebraic multiplicity at least $2$. If $T\in C_2(\mathbb{C}^5)$ lies in the Harnack part of $S$, then $0$ must have multiplicity $3$ (i.e., $\lambda=0$), and $T$ takes one of the following forms,
	\begin{equation}\label{T case12}
		T=\frac{2}{\sqrt{3}} \begin{bmatrix}
			0 & 1 & 0 & 0 & 0 \\
			0 & 0 & e^{i\theta} & 0 & 0 \\
			0 & 0 & 0 & e^{-i\theta} & 0 \\
			0 & 0 & 0 & 0 & 1 \\
			0 & 0 & 0 & 0 & 0
		\end{bmatrix},\quad \text{or} \quad 
		T=  \begin{bmatrix}
			0 & -\sqrt{3} & 0 & 0 & 0 \\
			0 & 0 & \frac{1}{\sqrt{2}}e^{i\theta} & 0 & 0 \\
			0 & 0 & 0 & \frac{1}{\sqrt{2}}e^{-i\theta} & 0 \\
			0 & 0 & 0 & 0 & -\sqrt{3} \\
			0 & 0 & 0 & 0 & 0
		\end{bmatrix},
	\end{equation} 
	with $\theta \in \mathbb{R}$.	
\end{proposition}

	\begin{equation}\label{eqfour}
		A \mathbf{R}(\overline{z}, \lambda) w(z) = -z^2 w^*(z) \mathbf{R}(z, \overline{\lambda}) B,
	\end{equation}

\begin{lemma}\label{lemma:Rn_unified}
	Let $R \in B(\mathbb{C}^3)$ such that $\sigma(R)=\{0, \lambda_1, \lambda_2\}$ with $\lambda_1, \lambda_2 \in \mathbb{D} \setminus \{0\}$. Then for any integer $n \geq 1$, the power $R^n$ can be expressed as
	\begin{equation}\label{eq:Rn_unified}
		R^n = A_n R^2 + B_n R,
	\end{equation}
	where the coefficients $A_n, B_n$ are given by
	\begin{equation}\label{eq:AnBn_unified}
		A_n = \begin{cases} 
			\displaystyle \frac{\lambda_1^{n-1} - \lambda_2^{n-1}}{\lambda_1 - \lambda_2}, & \text{if } \lambda_1 \neq \lambda_2, \\[10pt]
			(n-1)\lambda_1^{n-2}, & \text{if } \lambda_1 = \lambda_2,
		\end{cases}
		\quad 
		B_n = \begin{cases} 
			\displaystyle -\lambda_1 \lambda_2 \frac{\lambda_1^{n-2} - \lambda_2^{n-2}}{\lambda_1 - \lambda_2}, & \text{if } \lambda_1 \neq \lambda_2, \\[10pt]
			-(n-2)\lambda_1^{n-1}, & \text{if } \lambda_1 = \lambda_2.
		\end{cases}
	\end{equation}
	Consequently, for any $z \in \overline{\mathbb{D}}$, the resolvent expansion $(I-\overline{z}R)^{-1} = I_2 + c_1(z)R + c_2(z)R^2$ holds with coefficients
	\begin{align}
		c_1(z) &= \frac{\overline{z} - \overline{z}^2(\lambda_1 + \lambda_2)}{(1-\overline{z}\lambda_1)(1-\overline{z}\lambda_2)}, \label{c1} \\
		c_2(z) &= \frac{\overline{z}^2}{(1-\overline{z}\lambda_1)(1-\overline{z}\lambda_2)}. \label{c2}
	\end{align}
\end{lemma}

\begin{proposition}\label{Prop:case2}
	There exists no operator $T$ in the Harnack part of $S$ in $C_2(\mathbb{C}^{5})$ such that the block $R$ in the decomposition \eqref{form of T} satisfies $\sigma(R)=\{0, \lambda_1, \lambda_2\}$ with $\lambda_1, \lambda_2 \in \mathbb{D} \setminus \{0\}$ and $\lambda_1 \neq \lambda_2$.
\end{proposition}

\begin{proof}[Proof of Proposition \ref{Prop:case2}]
	Suppose, for contradiction, that there exists $T$ in $C_2(\mathbb{C}^{5})$ such that $T \stackrel{H}{\sim} S$ and $\sigma(R)=\{0, \lambda_1, \lambda_2\}$ with $\lambda_1, \lambda_2 \in \mathbb{D} \setminus \{0\}$. By Lemma~\ref{lemma:Rn_unified}, the resolvent $(I_2 - \overline{z}R)^{-1}$ expands as a quadratic polynomial in $R$. For $z \in \mathbb{D}$, we have
	\begin{equation}\label{resolvent_expansion}
		\mathbf{R}(\overline{z}, \lambda_1, \lambda_2) = (I_2 - \overline{z}R)^{-1} =  I_2 + c_1(z) R + c_2(z) R^2,
	\end{equation}
	where $c_1(z)$ and $c_2(z)$ are given by \eqref{c1} and \eqref{c2}. Note that these coefficients have poles at $z = 1/\overline{\lambda}_1$ and $z = 1/\overline{\lambda}_2$, which lie outside $\overline{\mathbb{D}}$ since $|\lambda_1|, |\lambda_2| < 1$.
	
	As in the proof of Proposition~\ref{Harnack_part_S_5}, the kernel condition $K_z^2(T)v(z) = 0$ yields the system
	\begin{equation}\label{system01_2} 
		\begin{cases}
			2v_0 + A \mathbf{R}(\overline{z}, \lambda_1, \lambda_2) w(z) - z^2 v_0 A \mathbf{R}(\overline{z}, \lambda_1, \lambda_2) B = 0,\\
			v_0 \mathbf{R}^*(z, \overline{\lambda}_1, \overline{\lambda}_2) A^* + K_z^2(R) w(z) - v_0 z^2 \mathbf{R}(\overline{z}, \lambda_1, \lambda_2) B = 0,\\
			v_0 B^* \mathbf{R}^*(z, \overline{\lambda}_1, \overline{\lambda}_2) A^* + B^* \mathbf{R}^*(z, \overline{\lambda}_1, \overline{\lambda}_2) w(z) - 2v_0 z^2 = 0.
		\end{cases} 
	\end{equation}
	Combining the first and third equations eliminates the constant term $2v_0$, yielding the analogue of \eqref{eqfour},
	\begin{equation}\label{eqfour_2}
		A \mathbf{R}(\overline{z}, \lambda_1, \lambda_2) w(z) = -z^2 w^*(z) \mathbf{R}(z, \overline{\lambda}_1, \overline{\lambda}_2) B.
	\end{equation}
	Let $P_1(\zeta) = (1 - \zeta\lambda_1)(1 - \zeta\lambda_2) = 1 - \sigma_1\zeta + \sigma_2\zeta^2$. The key equation \eqref{eqfour_2} can be rewritten as,
	\begin{align}\label{poly_identity}
		&A \left[ P_1(\overline{z}) I_2 + (\overline{z} - \overline{z}^2\sigma_1) R + \overline{z}^2 R^2 \right] w(z) \nonumber \\
		&\quad = -z^2 w^*(z)  \left[ P_1(\overline{z}) I_2 + (\overline{z} - \overline{z}^2\sigma_1) R + \overline{z}^2 R^2 \right] B.
	\end{align}
	Recall $w(z) = v_1 (1, 0, -z^2)^\top$. By analytic continuation, this identity holds for all $z \in \mathbb{C}$. Expanding both sides as polynomials in $z$ and comparing coefficients of $z^k$ for $k = 0, 1, \dots, 4$ yields, 
		\begin{equation}\label{coeff_z0}
			A (\sigma_2 I_2 - \sigma_1 R + R^2) \mathbf{e}_1 = \mathbf{e}_3^\top (\sigma_2 I_2 - \sigma_1 R + R^2) B,
		\end{equation} 
		\begin{equation}\label{coeff_z1}
			A (-\sigma_1 I_2 + R) \mathbf{e}_1 = \mathbf{e}_3^\top (-\sigma_1 I_2 + R) B,
		\end{equation}
		\begin{equation}
			 A \mathbf{e}_1 - A(-\sigma_1 R + R^2)\mathbf{e}_3
		 - \mathbf{e}_3^\top B + \mathbf{e}_1^\top (-\sigma_1 R + R^2) B = 0, \label{coeff_z2}
		\end{equation} 
	\begin{equation}
	  A(-\sigma_1 I_2 + R)\mathbf{e}_3 -  \mathbf{e}_1^\top (-\sigma_1 I_2 + R) B = 0, \label{coeff_z3}
		\end{equation} 
and
		\begin{equation}\label{coeff_z4}
			A \mathbf{e}_3 = \mathbf{e}_1^\top B.
		\end{equation}
	Equation \eqref{coeff_z4} immediately implies $a_3 = b_1$.
	
	Next, we analyze the first equation of \eqref{system01_2}. Let $Q(z) = (z-\lambda_1)(z-\lambda_2) = z^2 - \sigma_1 z + \sigma_2$. Multiplying by $Q(z)$ yields
	\begin{align}\label{eq:1c2}
		2v_0 Q(z) &+ A \left( Q(z)I_2 + (z-\sigma_1) R + R^{2} \right) w(z) \nonumber \\
		& - z^2 v_0 A \left( Q(z)I_2 + (z-\sigma_1) R + R^{2} \right) B = 0.
	\end{align}
	Evaluating at $z=0$ (so $Q(0)=\sigma_2$) gives
	\begin{equation}\label{eq:c21z0}
		A \left( \sigma_2 I_2 - \sigma_1 R + R^{2} \right) \mathbf{e}_1 = 2v_0 \sigma_2.
	\end{equation}
 Let $P(z) = (1-z\overline{\lambda}_1)(1-z\overline{\lambda}_2) = \overline{\sigma}_2 z^2 - \overline{\sigma}_1 z + 1$. The third equation of \eqref{system01_2} can be written as,
	\begin{align}\label{eq:3c2}
		v_0 B^* & \left( P(z)I_2 + (z-z^2\overline{\sigma}_1) R^* + z^2R^{*2} \right) B^* \nonumber \\
		& + B^* \left( P(z)I_2 + (z-z^2\overline{\sigma}_1) R^* + z^2R^{*2} \right) w(z)- 2z^2 v_0 P(z) = 0.
	\end{align}
The coefficients of $z^4$ gives
	\begin{equation}\label{eq:c21z4}
	v_1	B^* \left( \overline{\sigma}_2 I_2 - \overline{\sigma}_1 R^* + R^{*2} \right) \mathbf{e}_1 = -2v_0 \overline{\sigma}_2.
	\end{equation}

	We now turn to the second vector equation in \eqref{system01_2}. Writing first the resolvent terms explicitly,
	\begin{align}
		v_0\mathbf{R}^*(z, \overline{\lambda}_1, \overline{\lambda}_2)A^* &= v_0 \left( I_2 + \overline{c_1(z)} R^* + \overline{c_2(z)} R^{*2} \right) A^*, \label{app:term1} \\
		K_z^2(R) w(z) &= 2v_1 \begin{bmatrix} 1 \\ 0 \\ -z^2 \end{bmatrix} + v_1 \left( \overline{c_1(z)} R^* + c_1(z) R \right) \begin{bmatrix} 1 \\ 0 \\ -z^2 \end{bmatrix} \nonumber \\
		&\quad + v_1 \left( \overline{c_2(z)} R^{*2} + c_2(z) R^2 \right) \begin{bmatrix} 1 \\ 0 \\ -z^2 \end{bmatrix}, \label{app:term2} \\
		-v_0 z^2 \mathbf{R}(\overline{z}, \lambda_1, \lambda_2) B &= -v_0 z^2 B - v_0 z^2 c_1(z) R B - v_0 z^2 c_2(z) R^2 B. \label{app:term3}
	\end{align}
	Multiplying the sum of \eqref{app:term1}--\eqref{app:term3} by $P(z)\overline{P}(z)$, where
	\begin{align*}
		P(z) &= (1-z\overline{\lambda}_1)(1-z\overline{\lambda}_2) = \overline{\sigma}_2 z^2 - \overline{\sigma}_1 z + 1, \\
		\overline{P}(z) &= \overline{z}^2 (z-\lambda_1)(z-\lambda_2) = \overline{z}^2 Q(z),
	\end{align*}
we obtain then the following polynomial vector equation,
	\begin{align}\label{app:vector_cleared}
		v_0 &\Bigl( Q(z) P(z)  +  zQ(z) (1 - z\overline{\sigma}_1) R^*  +  z^2Q(z) R^{*2}\Bigr)   A^* \nonumber \\
		&+ v_1 \Bigl( Q(z) P(z)  +  zQ(z) (1 - z\overline{\sigma}_1) R^*  +  z^2Q(z) R^{*2}\Bigr)  \begin{bmatrix} 1 \\ 0 \\ -z^2 \end{bmatrix} \\
		&+ v_1 \Bigl( Q(z) P(z)+ P(z) (z - \overline{\sigma}_1) R +P(z) R^2 \Bigr)  \begin{bmatrix} 1 \\ 0 \\ -z^2 \end{bmatrix} \nonumber \\
		&\qquad \qquad - v_0 z^2\Bigl(  Q(z) P(z)  -  P(z) (z - \overline{\sigma}_1) R  - P(z) R^2\Bigr)   B = 0.\nonumber 
	\end{align}
	Extracting the coefficients of $z^0$ and $z^6$ gives, respectively,
	\begin{align}
		v_0 \sigma_2 A^* + v_1 \left( R^2 - \sigma_1 R + 2\sigma_2 I_2 \right) \mathbf{e}_1 &= 0, \label{eq2case2:z0} \\
		v_1 \left( R^{*2} - \overline{\sigma}_1 R^* + 2\overline{\sigma}_2 I_2 \right) \mathbf{e}_3 + v_0 \overline{\sigma}_2 B &= 0. \label{eq2case2:z6}
	\end{align}
	The coefficients of $z^2$ and $z^4$ yield the following vector equations,
	\begin{equation}\label{eq2case2:z2}   
		\begin{aligned}
			&v_0\Bigl( (1 + |\sigma_1|^2 + |\sigma_2|^2) +  (\sigma_1 + \sigma_2 \overline{\sigma}_1) R^*  + \sigma_2 R^{*2}\Bigr) A^* \\
			&+ v_1 \Bigl(\sigma_2 R^{*2} -(\sigma_1 + \sigma_2 \overline{\sigma}_1) R^* - (\overline{\sigma}_1 + \sigma_1 \overline{\sigma}_2) R + \overline{\sigma}_2 R^2 2v_1 (1 + |\sigma_1|^2 + |\sigma_2|^2)\Bigr) \mathbf{e}_1 \\
			&-v_1 \left( R^{2} - \overline{\sigma}_1 R + 2\sigma_2 I_2 \right) \mathbf{e}_3- v_0\Bigl( \sigma_2  - \sigma_1 R  + R^2\Bigr) B = 0,
		\end{aligned}
	\end{equation}
	and \begin{align}
		&	v_0 \Bigl( R^{*2} - \overline{\sigma}_1 R^* + \overline{\sigma}_2 I_2\Bigr) A^*+ v_1\Bigl( R^{*2} - \overline{\sigma}_1 R^* + 2\overline{\sigma}_2 I_2 \Bigr) \mathbf{e}_1 \nonumber \\
		&+ v_1\Bigl( -\sigma_2 R^{*2} + (\sigma_1 + \sigma_2 \overline{\sigma}_1) R^* + (\overline{\sigma}_1 + \sigma_1 \overline{\sigma}_2) R - \overline{\sigma}_2 R^2 - 2(1 + |\sigma_1|^2 + |\sigma_2|^2)\Bigr) \mathbf{e}_3  \label{eq2case2:z4}\\
		&- v_0 \Bigl(\overline{\sigma}_2 R^2 - (\overline{\sigma}_1 + \sigma_1 \overline{\sigma}_2) R + (1 + |\sigma_1|^2 + |\sigma_2|^2) I_2\Bigr) B = 0, \nonumber 
	\end{align}
	Taking the inner product of \eqref{eq2case2:z2} with $\mathbf{e}_1$ and adding it to the adjoint of \eqref{eq2case2:z4} applied to $\mathbf{e}_3$, then using  \eqref{coeff_z1}--\eqref{coeff_z3}, \eqref{eq2case2:z0}, and \eqref{eq2case2:z6}, we obtain after straightforward calculus 
	\begin{equation}\label{key_real2}
		v_0\left(1-\left|\sigma_2 \right|^2  \right) \left(\operatorname{Re}(a_1) -\operatorname{Re}(b_3)\right)-v_1\left( \sigma_1(\overline{r}_{11}+\overline{r}_{33})-\overline{\sigma}_1(r_{11}+r_{33})\right) = 0. 
	\end{equation}
	Since the second term must be purely imaginary, it vanishes, forcing
	\begin{equation}\label{key_real_eq}
		\operatorname{Re}(a_1) = \operatorname{Re}(b_3).	
	\end{equation}
	
	Rearranging \eqref{eq2case2:z0} and \eqref{eq2case2:z6} gives
	\begin{align}
		v_1 \left( R^2 - \sigma_1 R + \sigma_2 I_2 \right) \mathbf{e}_1 &= -v_0 \sigma_2 A^* - v_1\sigma_2 \mathbf{e}_1, \label{eqcase2:z0a} \\
		v_1 \left( R^{*2} - \overline{\sigma}_1 R^* + \overline{\sigma}_2 I_2 \right) \mathbf{e}_3 &= -v_0 \overline{\sigma}_2 B - v_1\overline{\sigma}_2 \mathbf{e}_3. \label{eqcase2:z0b}
	\end{align}
Thus
\begin{align*}
	v_1 A \left( R^2 - \sigma_1 R + \sigma_2 I_2 \right) \mathbf{e}_1 &= -v_0 \sigma_2 AA^* - v_1\sigma_2 A\mathbf{e}_1, \\
	v_1 B^*\left( R^{*2} - \overline{\sigma}_1 R^* + \overline{\sigma}_2 I_2 \right) \mathbf{e}_3 &= -v_0 \overline{\sigma}_2 B^*B - v_1\overline{\sigma}_2 B^*\mathbf{e}_3. 
\end{align*}
Hence by \eqref{eq:c21z0} and \eqref{eq:c21z0}, we get
\begin{align*}
v_0 \sigma_2 AA^* + v_1\sigma_2 A\mathbf{e}_1 &= 2v_1\sigma_2, \\
v_0 \overline{\sigma}_2 B^*B + v_1\overline{\sigma}_2 B^*\mathbf{e}_3 &=2v_1\overline{\sigma}_2 . 
\end{align*}
	Assuming $\sigma_2 = \lambda_1\lambda_2 \neq 0$, taking in account \eqref{key_real_eq}, then both $a_1$ and $b_3$ are real and satisfy the same quadratic equation
	\begin{equation}\label{quadratic_general}
		v_0 x^2 + v_1 x + C = 0,
	\end{equation}
	so,
	\begin{equation}\label{sum_roots}
		a_1 + b_3 = -\frac{v_1}{v_0}.
	\end{equation}
	On the other hand, direct spectral computation using trace properties shows  $\operatorname{Tr}(R^2 - \sigma_1 R) = -2\sigma_2$. Computing the trace via the standard basis and using \eqref{eqcase2:z0a}--\eqref{eqcase2:z0b} yields
	\begin{align*}
		v_1 \operatorname{Tr}(R^2 - \sigma_1 R) &=v_1\left\langle (R^2 - \sigma_1 R)\mathbf{e}_1, \mathbf{e}_1 \right\rangle+v_1\left\langle (R^2 - \sigma_1 R)\mathbf{e}_2, \mathbf{e}_2 \right\rangle+v_1\left\langle (R^2 - \sigma_1 R)\mathbf{e}_3, \mathbf{e}_3 \right\rangle \\
		&= \left\langle -v_0 \sigma_2 A^* - v_1\sigma_2 \mathbf{e}_1, \mathbf{e}_1 \right\rangle + v_1\left\langle (R^2 - \sigma_1 R)\mathbf{e}_2, \mathbf{e}_2 \right\rangle \\
		&\quad + \left\langle \mathbf{e}_3, -v_0 \overline{\sigma}_2 B - v_1 \overline{\sigma}_2 \mathbf{e}_3 \right\rangle \\
		&= -v_0(a_1+b_3)\sigma_2 - 4v_1 \sigma_2 + v_1\left\langle (R^2 - \sigma_1 R)\mathbf{e}_2, \mathbf{e}_2 \right\rangle \\
		&= -3v_1 \sigma_2 + v_1\left\langle (R^2 - \sigma_1 R)\mathbf{e}_2, \mathbf{e}_2 \right\rangle.
	\end{align*}
	Dividing by $v_1$ and comparing with $\operatorname{Tr}(R^2 - \sigma_1 R) = -2\sigma_2$ forces
	\[
	\left\langle (R^2 - \sigma_1 R)\mathbf{e}_2, \mathbf{e}_2 \right\rangle = \sigma_2.
	\]
	Substituting this back into the trace expansion of $R^2 - \sigma_1 R - \sigma_2 I_2$ gives
\begin{align*}
		\operatorname{Tr}(R^2 - \sigma_1 R - \sigma_2 I_2) &= \left\langle (R^2 - \sigma_1 R - \sigma_2 I_2) \mathbf{e}_1, \mathbf{e}_1 \right\rangle + v_1\left\langle(R^2 - \sigma_1 R - \sigma_2 I_2)\mathbf{e}_3, \mathbf{e}_3\right\rangle \\
		&=\left\langle (R^2 - \sigma_1 R)   \mathbf{e}_1, \mathbf{e}_1 \right\rangle + \left\langle(R^2 - \sigma_1 R)\mathbf{e}_3, \mathbf{e}_3\right\rangle-2 \sigma_2 \\
		&=-\frac{v_0}{v_1}(a_1+b_3)\sigma_2 - 5\sigma_2 = -4\sigma_2.
\end{align*}
	However, the characteristic equation implies $\operatorname{Tr}(R^2 - \sigma_1 R - \sigma_2 I_2) = -3\sigma_2$. Equating these expressions yields $-4\sigma_2 = -3\sigma_2$, which forces $\sigma_2 = 0$. This contradicts the assumption $\lambda_1, \lambda_2 \neq 0$. 
	
	Thus, at least one of $\lambda_1$ or $\lambda_2$ must be zero, reducing the problem to the case treated in Proposition~\ref{Harnack_part_S_5}, where we established $\lambda_1 = \lambda_2 = 0$. Consequently, no operator $T$ with $\sigma(R)=\{0, \lambda_1, \lambda_2\}$ and $\lambda_1, \lambda_2 \neq 0$ can belong to the Harnack part of $S$. This completes the proof.
\end{proof}

\end{document}
